\section{The wider investigation, by method}\label{sec:atlas}
The case studies above give explicit constructions and proofs for the
principal unfinished routes as well as the selected obstructions.
This shorter catalog retains related variants, exact rediscoveries and
bounded experiments. An entry described only here retains its archived
scope; it is not promoted to a theorem proved in this supplement.
The linked labels lead to the source register. They let a reader pursue
an exploratory variant without making the research logs prerequisites
for the case studies.

\subsection{Rational Euler and two-endpoint Hermite sources}
These constructions cancel Taylor coefficients at one or two points.
The aim is to force a small remainder while keeping its coefficients
arithmetically manageable. Solving the interpolation equations is only
the first step; their endpoint values still need nonvanishing and
primitive-size estimates.

Records \recordid{H01}--\recordid{H03} extend~\eqref{eq:euler-exact} to two-endpoint sources.
Section~\ref{sec:endpoint-filter-case} proves the four-by-five reduction,
including singular cases and the constant coefficient. The circular
Hermite specialization has an archived eventual rank result and a signed
error asymptotic proportional to
$(-1)^n n^{5/2}\log^2 n\,((\sqrt2-1)/2)^n$.
Its actual primitive height remains uncontrolled. Factorial rescaling
of the entire pair does not change that question.

One arithmetic ingredient, \recordid{H02}, can be stated independently of that
source. Products of $r$ distinct integers in $[1,N]$, $N\geq r$,
are all divisors of
\[
 W_r(N)=\prod_{j=1}^r\operatorname{lcm}(1,\ldots,\lfloor N/j\rfloor).
\]
At a prime $p$ the maximum possible product valuation is
$\sum_{k\geq1}\min(r,\lfloor N/p^k\rfloor)$; this is exactly
$v_p(W_r(N))$. The maximum is attained by choosing the $r$ integers
with highest valuations, since the sets of multiples of $p^k$ are nested.
For the strict integration indices of the archived source this gives
a sufficient coefficient denominator. In particular
$\log W_4(N)\sim(1+1/2+1/3+1/4)N=25N/12$ by the prime number
theorem. It is not an endpoint primitive-value bound.

\subsection{Continued fractions, derivatives and asymptotic summation}
These methods improve a series tail by a rational approximation or an
asymptotic correction. Differentiation can change the target constant,
but also changes the denominator cost. It is essential to specify which
endpoint and truncation parameters are allowed to grow together.

\recordid{S04} and \recordid{S05} are Theorems~\ref{thm:boole} and~\ref{thm:cf}; their
different quantifiers matter. Euler--Boole covers all endpoint and
truncation schedules, whereas the continued-fraction theorem fixes the
endpoint and restricts arithmetic to even depth. \recordid{E05} concerns rational
gauges (invertible rational changes of dependent variables) of a
particular two-component tail system: the archived argument
excludes its desired polynomial unimodular transfer. It says nothing
about factorial gauges, higher-rank systems or unrelated recurrences.
The binomial-shift reduction at the end of Section~\ref{sec:cf}
also prevents counting an old scalar source as a new continued-fraction
construction.

\subsection{Modular and Eichler sources}
Modular transformations relate expansions of a function at different
points. The hope is to combine arithmetic information from those
expansions with an analytic estimate improved by the symmetry.
The experiments below test cancellation of unwanted constants, removal
of selected prime denominators, and the singularities that remain.
The Eichler construction integrates a modular form; its transformation
law introduces a polynomial containing auxiliary constants, which
must be eliminated before isolating $\theta$.

Section~\ref{sec:modular-case} develops the period cancellation of
\recordid{M01}, a prime-deleted cusp source, and its fixed-filter natural
boundary. \recordid{M02} also constructs quadratic prime filters
compatible with the Fricke transformation, obtaining selected savings
in both endpoint germs. Neither transformation identities nor local
savings alone establish favorable primitive arithmetic for growing
filters. \recordid{M03} records uncancelled interior poles in proposed higher-level
coordinates and prefactors. A later annihilator removes ordinary and
shifted pole families but leaves six fixed-point poles. These are
limitations of those coordinates and raw coefficient estimates.

\recordid{M04} compares three specified rational-system templates of ranks
$9,17,18$. Its optimized displayed logarithmic derivative bound is
below $1.3$, against a required threshold above $4$ in that test.
This closes the prescribed comparison, not all modular coordinates.
The exact template and source choices are retained in the record;
they are essential hypotheses, not an incidental implementation detail.

\subsection{Polylogarithm ladders and functional equations}
Functional identities replace one special value by combinations at
other arguments. Small arguments improve series convergence, and shared
arguments can permit joint cancellation. The remaining task is to
control all coefficients and auxiliary logarithms in the resulting form.

\recordid{L01} is the completed Broadhurst identity used in
Section~\ref{sec:mellin}; its attribution is part of the statement.
\recordid{L02} studies equal-speed scalar Taylor truncations of one rational
realization. Using the $c_m,z_m$ from Section~\ref{sec:mellin}, set
$\ell(t)=-\log(1-t/2)$ and, for real $t$ near zero,
\[
 H(t)=\sum_{m\in\{2,3,5\}}c_m\Im\left(
 \Li_4(z_mt)+\frac{m\ell(t)}2\Li_3(z_mt)
                  +\frac{m^2\ell(t)^2}{12}\Li_2(z_mt)\right).
\]
Write $H(t)=\sum_{j\geq1}h_jt^j$; its coefficients are rational
and $H(1)=\theta$. Let $q_N$ be the reduced denominator of
$\sum_{j=1}^N h_j$. The archived exact
dyadic exponents are
\[
 v_2(q_{4m})=10m+4,\quad v_2(q_{4m+1})=10m+5,\quad
 v_2(q_{4m+2})=10m+11,\quad v_2(q_{4m+3})=10m+12,
\]
with the first formula restricted to $m\geq1$. Combined with the
signed error estimate, these give a primitive obstruction for that
scalar truncation, not for other uses of the ladder.

The common-base construction \recordid{L03} and its precise missing
odd-prime estimate are developed in Section~\ref{sec:common-base}.
The compatible pole multipliers \recordid{L04} and their eventual signed
error asymptotic are developed in Section~\ref{sec:compatible-case}.
Neither is excluded by the scalar Taylor-truncation result above.
\recordid{L05} specializes Gangl's weight-four functional equation~\cite{Gangl}
to a 21-argument algebraic presentation. The archive checks the symbol
identity and tracks completion terms. A naive trace after separate
inversions fails; no theorem of numerical period independence follows
from a symbol computation.

\subsection{Slow dilations, prime deletion and sparse Pad\'e}
These attempts modify the generating series to reduce its arithmetic
cost: spread out its exponents, remove contributions from selected
primes, or retain a sparse set of terms. Such changes also alter the
singularities and the available approximation estimates, so an arithmetic
saving alone does not settle the construction.

Section~\ref{sec:compatible-case} defines the slow-dilation germ and
proves its regular-word coefficient bound. The sharper archived result
\recordid{A01}, for fixed odd $k$, finds coefficient
denominator rates between $3/2+1/(k+2)$ and
$3/2+1/k+1/(2k+4)$ in the normalization defined in its record.
These are fixed-$k$ coefficient statements, not primitive determinant
rates or estimates uniform in a jointly growing $k$.
\recordid{A02} relates a bordered determinant's content to maximal-minor gcds
and modular Pad\'e defects. Here $H$ is the archived rational moment
matrix, $b$ its bordering vector, and $D$ clears their entries.
With $T=DH$, $u=Db$,
$d=\det T$, $m=\operatorname{adj}(T)u$, and
$Q=u^T\operatorname{adj}(T)u$, the specialized gcds obey
$\Delta\mid g\mid6D\Delta^2$, where
$g=\gcd(6Dd,Q)$ and $\Delta=\gcd(d,m_i)$.
An example contains prime $23$ absent from the coefficient denominator.
The determinant identities are classical; the bounds do not settle
all asymptotic primitive content.

\recordid{A03} concerns failed fixed-cut absolute-contour estimates for $k=3,5$.
It is not a lower bound on actual determinant size. \recordid{A06} explains why
a proposed same-variable positivity repair fails: the slow filter has
genuine nonreal branch points satisfying
$z^k=-(\sqrt2-1)^2$, inconsistent with the proposed compact
real-supported Cauchy transform after the analyzed algebraic repairs:
such a transform must be holomorphic off its real support.
Different changes of variable are outside that statement.

Section~\ref{sec:endpoint-filter-case} proves the selected-prime saving
of \recordid{A04} and the sparse determinant collapse of \recordid{A05}.
The latter's further primitive-growth theorem is restricted to an
oversieved regime: its cutoff prime is at least the fourth power of
its size parameter. The natural cutoff was not excluded.

\subsection{Positive determinants and interleaved frames}\label{sec:interleaved}
Changing the row and column nodes of a determinant may improve the
balance between its analytic size and its prime denominators. The key
distinction is between changing a basis within the same polynomial
pencil, which preserves its primitive polynomial up to sign, and choosing a
different set of rows, columns or moments, which can change the problem.

\recordid{C01}--\recordid{C08} are the companion paper's positive cubic-tail kernel, uniform
Cauchy comparison, consecutive-principal obstruction, difference-Gram
construction and exact large-prime primitive arithmetic. Its two
primitive lower exponents are $2-2\log2$ and $6-\log4$ on the $n^2$
scale. \recordid{S08} is a separate archived obstruction for an overshifted dense
grid with $m\geq768r^4$; it does not address proportional shifts.
The unrestricted source cone is covered in Section~\ref{sec:cone}.

Section~\ref{sec:interleaved-case} develops the nonprincipal experiment
\recordid{P06}, including its local prime layers and the distinction
between integer clearing and rational primitive normalization.
Section~\ref{sec:variation} develops the separate moment-differentiation
construction \recordid{P07}. Neither remaining arithmetic question is
settled by the companion paper's principal-family theorems.

\subsection{Hypergeometric and quarter-pole constructions}
Factorial ratios and hypergeometric identities can provide rational
coefficients together with exact remainder formulas. Shifting or
differentiating their parameters changes the special values that occur.
The tests here seek a combination containing only $\theta$ and a
rational term, then examine the denominator powers and common factors.

\recordid{Q01} differentiates an integer-pole tail to isolate $\theta$, with
signed error proportional to
\[
 (-1)^{n+1}n^{-5/2}\log^2n(27/256)^n.
\]
Its displayed clearing multiplier retains fifth odd-prime layers;
the resulting unfavorable bound is not a theorem about full primitive
content. \recordid{Q02} tests three nonsymmetric half-parameter jets and a
120-element Thomae orbit at one fixture. Its six primitive eliminations
are finite evidence, not a classification of transformations.

\recordid{Q03} supplies an all-index moved-base Whipple identity, a positive
terminating coefficient expression, dyadic valuations and eventual
signed nonvanishing. Its proved multipliers obey a divisibility relation,
and integer clearing of the selected cofactor fails even before
odd-prime costs. Unbounded odd common content remains outside the result.
The symmetric candidate was exactly an earlier archived family.

Section~\ref{sec:paired-case} gives the complete quarter-shifted paired
construction \recordid{Q04}, the proof mechanism and standalone algebraic
certificates. It also specifies and proves the transferred-estimate
limitation \recordid{Q05}. The earlier numerator mismatch was repaired;
it is not a reason to reject the corrected identities. The unresolved
questions concern projected nonvanishing and primitive size.

\subsection{Coupled integrals, geometry, and exact rediscoveries}
Coupling several integrals or introducing extra parameters can cancel
unwanted constants while retaining a positive expression. One must then
check both the resulting arithmetic and whether the apparent new
construction is actually a reparametrization of an earlier one.

\recordid{P01} and \recordid{O01} navigate the earlier lower-period cancellations, coupled
holonomy, lattice, radial and orthogonal constructions. They are umbrella
records, not single consolidated theorems. Some fitted recurrences were
conjectural before later certification; those stages remain distinct.

\recordid{P02} identifies the positive dimensional family exactly as
$J_n=(n!)^4 4^{-n}L_{n+1,n+1}$ in the archived balanced-ladder
notation. Here $J_n$ is the proposed coupled integral and $L_{b,h}$
the earlier family with two integer parameters; their full integrands
are given in the record. The displayed factor is a rational scalar,
so primitive normalization removes it. This is the precise sense in
which the proposed family is a rediscovery. Section~\ref{subsec:five-period} gives the full rational-series
construction and five-period reduction of \recordid{P03}, followed by
an explicit cofactor elimination. Section~\ref{sec:quarter-case} develops
the squared and unsquared quarter-pole pencils \recordid{P04}--\recordid{P05},
including the latter's conditional nonvanishing proof. Their actual
primitive asymptotics remain unresolved.

\recordid{S07} is a quarter-Gamma fourth-cumulant extraction from the weighted
projective data $(1,3,4,4,4,4,4)$. Its scalar sequence uses finite
Hurwitz fourth-power sums, with weights $1,3,4,4,4,4,4$ and total $24$,
at the two quarter arguments. The fourth logarithmic coefficient is
formed with all lower products retained. Its archived arithmetic proof
gives reduced-denominator logarithm at least $24n+o(n)$, whereas
the error is of order $n^{-5}$. That scalar obstruction does not certify
the larger geometric interpretation or exclude the separate nonconstant
deformation. No geometric theorem is needed for the finite-sum obstruction.

\recordid{E09} records other exact collapses: signed divisibility/LCM Gram
determinants become linear approximants by rank one, and complete Boolean
bases reproduce an Euler product. Polynomial Dirichlet convolution retains
the fourth denominator layer but also pays the hypothetical rational
denominator to the polynomial degree. Untested bases are not excluded.

\subsection{Euler functionals and proposed analytic bridges}
Here the main issue is whether a proposed connection between two kinds
of information is valid: for example, whether a $p$-adic functional
estimate controls the relevant coefficients, or whether an analytic
property supplies the arithmetic conclusion needed for irrationality.
A counterexample to such a connection closes only the stated implication.

\recordid{E01} computes a norm on polynomials taking integral values on
$\Z_p$, rather than on polynomials with integral monomial coefficients.
Define the Euler functional by
$\mathcal E(x^j)=E_j(0)$, where
$\sum_{j\geq0}E_j(0)z^j/j!=2/(e^z+1)$, and set
$\tau(P)=-\mathcal E(P''')/12$.
For $p\geq5$, let $m_p(h)$ be the least degree for which an
integer-valued polynomial $P$ satisfies $v_p(\tau(P))=-h$.
The exact onset formula is
$m_p(3q+r)=(3-r)p^q+rp^{q+1}$ for integers $q\geq0$, $0\leq r<3$;
the first three positive layers start at $p+2,2p+1,3p$.
This says nothing by itself
about rational poles or primitive determinant content.

\recordid{E02} applies Euler distribution to the whole polynomial-plus-pole
functional, with the affine target change
$Y=\chi_4(p)p^4X+C_p$. The positive csch functional was already
present in Carreira~\cite[Section 11]{Carreira}. Neither the distribution
identity nor a formal target change equates real and $p$-adic limits.
\recordid{E03} supplies countermodels to particular proposed analytic implications,
including an E-function Stokes inference, a compact-support inference
from rational moments, and a spectral-raising bridge. Each countermodel
has its stated hypotheses; none excludes the entire surrounding subject.

\subsection{Cyclotomic parity and rationality-dependent constructions}
The remaining tests use algebraic conjugation or begin by supposing
$\theta$ rational. In the first case a trace can vanish by symmetry;
in the second, an integrality claim or a differential construction must
actually distinguish the hypothetical rational value. Neither property
follows merely from writing down the corresponding formal expression.

The direct even-power signature extension \recordid{E08} lies in a cyclotomic
component negated by an involution fixing $\Q$, so its prescribed
rational traces vanish. This parity statement does not apply to the
nonlinear norm of Section~\ref{sec:cyclotomic}; that family needed its
own all-conjugate argument.

Section~\ref{subsec:niven-case} constructs the rationality-dependent
Niven source \recordid{E07} and proves the surviving origin denominator,
endpoint derivative and differential-field obstructions. Its scope is
the specified integrality and coordinate proposals.

\recordid{E06} subtracts the prescribed polar terms of a rationality-dependent
$q$-series. At $q=(p-1)/p$, its $p^2$ truncation has exact denominator
$p^{p^2}$ and growing numerator. This is a result for that rational
sampling and truncation, not for all $q$-series. \recordid{E04} is the different
monodromy theorem proved in Section~\ref{sec:mono}. In both cases the
hypothesis of rationality supplies no unspecified extra geometric or
$p$-adic relation.
